%%=============================================================================
%% Probleemanalyse en context
%%=============================================================================

\chapter{\IfLanguageName{dutch}{Probleemanalyse en context}{Problem analysis and context}}%
\label{ch:probleemanalyse}

Dit hoofdstuk werkt de probleemstelling verder uit door het toepassingsdomein van het onderzoek
, het PLAPO-lijnplanningsplatform binnen ArcelorMittal, gedetailleerd in
kaart te brengen. De uitgewerkte context vormt het uitgangspunt voor zowel de reverse engineering in
Hoofdstuk~\ref{ch:reverse-engineering} als de proof-of-concept in Hoofdstuk~\ref{ch:poc}.

\section{Het PLAPO-platform}
\label{sec:plapo}

PLAPO is de interne benaming voor het lijnplanningsplatform van ArcelorMittal. Het platform vertaalt
binnenkomende productieaanvragen naar concrete planningssegmenten per productiestap. Een segment beschrijft
welke materialen, door een specifieke installatie verwerkt kunnen worden.

Architecturaal bestaat PLAPO uit twee soorten modules:
\begin{description}
    \item[Generieke modules] zijn verantwoordelijk voor de afhandeling van een planningsaanvraag,
        het opbouwen van de gemeenschappelijke datacontext en het routeren naar de juiste specifieke
        module. Deze generieke modules zijn reeds volledig in PL/I ge\"implementeerd.
    \item[Specifieke modules] bevatten de domeinspecifieke logica per installatie of proces.
        Voorbeelden zijn de continue gietlijnen, de warmwalserij en de Batch Annealing Facility~(BAF).
        Op \'e\'en uitzondering na zijn alle specifieke modules eveneens reeds gemigreerd naar PL/I; alleen
        de BAF-module draait nog in CA~Gen.
\end{description}

Figuur~\ref{fig:Plapo} toont de positie van de BAF-module binnen het bredere PLAPO-landschap. De
generieke modules sturen, op basis van de aanvraagparameters, de juiste specifieke module aan; deze
laatste schrijft vervolgens de planningssegmenten weg in een sequenti\"eel uitvoerbestand dat door de
planningstool wordt ingelezen.

\section{De Batch Annealing Facility (BAF)}
\label{sec:baf}

Batch annealing is een warmtebehandeling waarbij meerdere rollen ijzer onder een vacu\"umkap geplaatst en
gecontroleerd verwarmd worden, met als doel de mechanische eigenschappen (hardheid, vervormbaarheid) te
verbeteren. Een BAF-installatie bestaat uit een aantal werkvelden waarop telkens een rolstapel kan worden
geplaatst. De planning bestaat eruit om aan de hand van enkele parameters een lijst van de planbare rollen op 
te stellen.

Concreet moet de BAF-planningsmodule per planbaar materiaal:
\begin{enumerate}
    \item Materiaaleigenschappen (afmetingen, kwaliteit, bestemming) bepalen.
    \item Productie-orders koppelen.
    \item Op basis van een set parameters en regels een of meer planningssegmenten samenstellen.
    \item Deze segmenten naar het uitvoerbestand schrijven in een vooraf afgesproken layout.
\end{enumerate}

\section{De CA~Gen-module TPPBAMT}
\label{sec:tppbamt}

Binnen het BAF-luik vormt de CA~Gen-module \texttt{TPPBAMT} het hart van de specifieke materiaallogica. 
Deze module genereert de specifieke materiaal-segmenten van alle planbare materialen voor de BAF-planning. 
Vanuit de generieke modules wordt \texttt{TPPBAMT} aangeroepen met een inputparamters en maakt de nodige
materiaalsegmenten en schrijft deze weg naar het output-bestand.

\section{Bedrijfsdoelstellingen van de migratie}
\label{sec:doelstellingen-migratie}

Het inhoudelijk doel van het bredere migratietraject (waarvan deze bachelorproef de
methodologische basis legt) is drieledig:
\begin{enumerate}
    \item De afhankelijkheid van CA~Gen-licenties en CA~Gen-expertise be\"eindigen.
    \item Het PLAPO-platform consolideren in \'e\'en doeltaal (PL/I), wat onderhoud en doorontwikkeling
        vereenvoudigt.
    \item De impliciete migratiekennis die in eerdere modulemigraties is opgebouwd, expliciet en
        herbruikbaar maken.
\end{enumerate}
Dit derde doel is precies wat het in Hoofdstuk~\ref{ch:framework} ontworpen framework realiseert.
